\section{Experiment Index}
\label{sec:experiment-index}

The code is the proof of last resort. This appendix indexes the key experiments by topic. The full numbered
series, p1 through p200 and beyond, is the running log. The checks below are the load-bearing ones cited in the
paper. The complete record is in \cite{pollard2026vibetest}.

\subsection{Substrate and builders}

\begin{itemize}
  \item The cell-direct builder holds the hyperbolic cell-graph builder and the flat-lattice builder for
  Euclidean controls.
  \item The Coxeter growth module holds the Steinberg growth series and the shell growth ratios.
  \item The Schl\"afli parsing and Coxeter machinery handle symbol reading and the reflection engine.
\end{itemize}

\subsection{The rule, conservation, and reversibility}

\begin{itemize}
  \item The pure-rule verification, with the parallel and serial runs agreeing exactly at the
  multi-million-cell scale.
  \item The reversibility check, running the rule forward and then backward and recovering the start state with
  zero error.
  \item The charge-conservation check across the rule and its conserving variants, with the total tone held
  exactly fixed.
  \item The determinism-versus-stochastic check, which shows a stochastic rule loses the light cone while the
  deterministic rule keeps it, the source of momentum.
  \item The reversible conserving rule-variant sweep, mapping which rule families show faithful emergence and
  which stay sterile or chaotic.
\end{itemize}

\subsection{Tessellation sweeps}

These produce the per-dimension data in Appendix~\ref{sec:tessellation-catalog}.

\begin{itemize}
  \item The per-dimension tessellation sweeps, one each for 2D, 3D, 4D, and 5D.
  \item The cross-dimension ranking battery.
\end{itemize}

\subsection{Spin, fermions, matter}

\begin{itemize}
  \item The Skyrme-twist stabilizer measured by the KPM sea energy, which is positive.
  \item The 3D soliton check.
  \item The p190 and p206 series, spin-statistics and the hopfion fermion.
\end{itemize}

\subsection{Persistence, defects, and the negatives}

\begin{itemize}
  \item The churn baseline, p101, where the pure rule from generic data gives structured churn with no
  persistent localized self, the load-bearing negative the rest of the program must beat.
  \item The persistent-pattern and glider search with identity-through-turnover scoring, on which a Conway Life
  glider on the cubic cusp persists cleanly while a generic blob disperses.
  \item The topological-protection checks, the conservation of the winding or Hopf charge, and the pair
  creation and annihilation of opposite defects.
\end{itemize}

\subsection{Dirac, Lorentz, isotropy}

\begin{itemize}
  \item The light cone and $D_4$ isotropy checks.
  \item p193, p211, p230 for dispersion, and p124, p125, p150, p154 for Lorentz, boost, and wave isotropy.
\end{itemize}

\subsection{Gauge and Standard Model}

\begin{itemize}
  \item p221, p225, p227, p228, p231, p234, covering $\mathfrak{so}(10)$, breaking,
  $\sin^2\theta_W$, and the Wilson loop.
  \item The renormalization-group checks for the RG elevation, which produces predictions.
\end{itemize}

\subsection{Gravity and holography}

\begin{itemize}
  \item The bulk correlator by Bethe recursion and tensor network, giving $1/r^2$.
  \item p224, p239, p218, the cusp Newton $1/r$ by the difference method, and holography.
\end{itemize}

\subsection{Cosmology and expansion}

\begin{itemize}
  \item The cosmology and anisotropy check, the expansion law $a(t) = e^{Ht}$ and the cubic anisotropy bound.
  \item The p10 and p13 series, the cosmological constant and the growth-equals-expansion result.
\end{itemize}

\subsection{Renormalization, the coarse-graining tower, and rarity}

\begin{itemize}
  \item The persistence-by-scale check and the slow conserved-density mode, a long-lived coarse field emerging
  from the churn.
  \item The radial coarse-graining tower, p187, a negative, radial blocking gives no hierarchical persistence
  tower.
  \item The one-dimensional decimation check, where block-spin summation matches the analytic recursion
  $K' = \tfrac{1}{2}\ln\cosh 2K$ exactly and flows to the Gaussian fixed point.
  \item The rarity cascade, p181, the multiplicative per-rung charge fraction that reproduces the cosmic
  one-part-in-a-hundred-trillion order, with the supporting measures, p183, the integration spectrum, the
  thin-film dimension, and the condensation density threshold.
\end{itemize}

\subsection{Physical space, the cusp}

\begin{itemize}
  \item The cusp-convergence, emergent-space, horosphere-reality, physics-on-real-space, and bulk-cusp-interface
  checks, the cusp-is-3D-space measurements and the $SO(4) \to SO(3)$ projection.
  \item p142 and p144, horosphere flatness and dynamics.
\end{itemize}

\subsection{Computation and universality}

\begin{itemize}
  \item The information-transport check, signals moving at the light-cone speed.
  \item The ternary signed-majority gate computing NAND, and the elementary rule $110$ built from it.
  \item The reversible charge-conserving register machine, with registers held as charge in addressed
  subtrees.
  \item Conway's Game of Life run on the cubic cusp, a strongly universal cellular automaton.
  \item The solved $\{3,4,3,4\}$ neighbour automaton and the $O(\log n)$ self-describing addressing.
\end{itemize}

\subsection{The emergent ladder and observers}

\begin{itemize}
  \item Causal emergence, p185, the coarse macro acting as a stronger cause than the micro, a measured
  downward causation.
  \item The coarse-grained self search redone on the pure rule with no added cohesion, a recorded negative
  consistent with p101.
  \item Reproduction, p120, self-copying patterns, with the recorded negative being none on the bare rule.
  \item The integration measure that picks out selves, the strange-loop self-model, and the global-workspace
  hub used in the mind and consciousness chapters.
  \item These ladder checks are toy-scale, and the upper rungs are the most speculative tier of the program.
\end{itemize}

\subsection{Cosmology, the open question}

\begin{itemize}
  \item The eternal-bootstrap check, the reversible eternal-churn viability test.
\end{itemize}

\subsection{Note on the log}

The p-numbered checks are an exploratory log. Not all are cited, and some early ones were superseded. The naive
finite-patch gravity, for example, was replaced by the difference method. The checks named explicitly above are the
current, clean results. The whole computational record is in \cite{pollard2026vibetest}.
